% Generated draft: requires mathematical and source review.
\documentclass[11pt]{article}
\usepackage[margin=25mm]{geometry}
\usepackage{fontspec}
\setmainfont{Latin Modern Roman}
\usepackage{amsmath,amssymb,mathtools,unicode-math}
\setmathfont{Latin Modern Math}
\usepackage{enumitem,needspace,xurl,hyperref}
\setlength{\parindent}{0pt}
\setlength{\parskip}{5pt}
\setlength{\emergencystretch}{3em}
\newcommand{\N}{\mathbb{N}}
\newcommand{\Z}{\mathbb{Z}}
\newcommand{\Q}{\mathbb{Q}}
\newcommand{\R}{\mathbb{R}}
\newcommand{\C}{\mathbb{C}}
\newcommand{\F}{\mathbb{F}}
\newcommand{\A}{\mathbb{A}}
\newcommand{\PP}{\mathbb{P}}
\newcommand{\E}{\mathbb{E}}
\newcommand{\eps}{\varepsilon}
\newcommand{\dd}{\,\mathrm{d}}
\begin{document}
\section*{UnsolvedMath ID: 20002472}
\textbf{Source problem code:} catalogue entry
\subsection{Definitions and mathematical statement}
Put $\mathbb T=\R/\Z$. John Mather's question asks whether there exist a $C^3$ twist map $f:\mathbb T\times\R\to\mathbb T\times\R$ and an invariant circle $\Gamma$ such that $f|_\Gamma$ has irrational rotation number but is not topologically conjugate to a rigid rotation. A conjugacy means a homeomorphism $q:\Gamma\to\mathbb T$ with $q\circ f|_\Gamma\circ q^{-1}(\theta)=\theta+\alpha$ for some $\alpha\notin\Q/\Z$. For an orientation-preserving circle homeomorphism with lift $F:\R\to\R$, its rotation number is $\lim_{n\to\infty}(F^n(x)-x)/n$ modulo $\Z$. The source recalls the existence of $C^1$ Denjoy examples on the circle and asks whether their behavior can occur on invariant curves of $C^3$ twist maps.

The source does not state its exact twist-map hypotheses or regularity of $\Gamma$. In the usual monotone-twist convention a lift $(\Theta,R)$ has $\partial_r\Theta>0$; area preservation, exactness, essentiality of the curve, and the sign convention must be checked against the intended setting. The aid's reduction applies to invariant graphs $\Gamma=\{(\theta,\gamma(\theta))\}$ with an induced orientation-preserving irrational circle homeomorphism $h$.

Its reported full-support criterion is
\[
 h\text{ is conjugate to }R_\alpha
 \ \Longleftrightarrow\ \exists\text{ invariant Borel probability }\mu
 \text{ with }\operatorname{supp}\mu=\mathbb T.
\]
Full support means every nonempty open arc has positive measure. Thus any counterexample in this graph setting must have a proper Cantor support $K$ for its invariant probability, and every component arc $I$ of $\mathbb T\setminus K$ is wandering: $h^n(I)$ are pairwise disjoint for distinct integers $n$. The whole invariant graph then strictly contains the measure support; in particular it cannot itself be the support of an invariant minimizing probability in the relevant variational twist-map setting.

The aid emphasizes two limitations. First, every irrational orientation-preserving circle homeomorphism is chain recurrent at every point, including Denjoy examples. Here chain recurrence means that for every $\eps>0$ and every $x$ there is an $\eps$-pseudo-orbit $x=x_0,\ldots,x_N=x$, $N\ge1$, satisfying $d(h(x_j),x_{j+1})<\eps$. It therefore does not certify conjugacy. Second, a $C^2$ invariant graph, giving a $C^2$ induced circle diffeomorphism in this setting, is excluded as a counterexample by Denjoy's theorem. None of these reductions constructs a $C^3$ twist-map example or proves that none exists.
\subsection{Short English statement}
The question is whether a smooth ambient twist map can carry an irrational invariant circle with Denjoy dynamics. The supplied reduction says such a circle must contain wandering arcs outside its invariant measure's Cantor support. Full-support invariant measure excludes the pathology; chain recurrence alone does not.
\subsection{Sources}
\begin{itemize}
\item Local UnsolvedMath registry, permanent ID 20002472, source code catalogue entry; created and updated 14 August 2026. Canonical \texttt{aim-pdes-notes.json}, zero-based \texttt{notes[83]}, tag \texttt{problem}, item 9 attributed to John Mather.
\item AIM ``New connections between dynamical systems and PDE's'': \url{https://aimath.org/WWN/dynpde/dynpde.pdf}. The supplied verification identifies the canonical item and marks its wording exact. The terminal 1 is extraction pagination. Original remarks and literature field are empty.
\item Research attempt 2: \texttt{partial\_result}, problem status at run open, rigor label \texttt{proved}, novelty confidence low. Denjoy's example and theorem are named, but no complete authoritative bibliographic references or full twist-map hypotheses are supplied.
\item Ulam AI / UnsolvedMath Contributors, CC BY 4.0. Machine-generated research aid, not peer reviewed; independently verify before citation.
\end{itemize}
\end{document}
